% Report template. `make pdf` copies this file into build/ and compiles it there with pdflatex,
% so the \input and \includegraphics paths below are relative to build/.
\documentclass[11pt]{article}

\usepackage[utf8]{inputenc}
\usepackage[margin=1in]{geometry}
\usepackage{amsmath,amssymb}
\usepackage{booktabs}
\usepackage{graphicx}
\usepackage{float}
\usepackage[colorlinks=true,linkcolor=blue!50!black,urlcolor=blue!50!black]{hyperref}
\usepackage{xcolor}

\graphicspath{{figures/}}
\setlength{\parindent}{0pt}
\setlength{\parskip}{0.6em}
\sloppy
\emergencystretch=2em

\title{Numerical and symbolic models for physics results\\in the OpenAI mathematics release\\[0.2em]\large Fork notes generated by \texttt{tools/}}
\date{\today}
\author{}

\begin{document}
\maketitle

\begin{abstract}
This document accompanies the scripts in \texttt{tools/}. For four families of physics-related results in the repository
(Bloch's law for the quantum Heisenberg ferromagnet, mass scaling in two-dimensional $\mathrm{O}(N)$ models, the XY model's BKT
scaling, and the spin-1 Haldane gap) it builds a small model, computes with numpy, scipy, sympy and matplotlib, and checks the
statements of the repository against those computations where that is possible. The models are illustrations: none of them reproduces
or re-proves a repository theorem, and each section lists what is and is not tested.
\end{abstract}

\section{Overview}

The repository's entries are cited by their number in \texttt{CONTENTS.md}. The sections below correspond to:
entry 271 (Bloch's law), entry 215 (two-dimensional $\mathrm{O}(N)$ models), entry 216 (XY model) and entry 268 (Haldane gap).
Section~\ref{sec:chair44} is different in kind: it is a constraint model built from the contact atlas of an external preprint on a seven-cube aperiodic tile, not from a repository entry. The four entries above were chosen because each admits an independent, checkable computation; the repository's Yang--Mills discussion is
in \texttt{YANGMILLS.md}, and the repository contains no four-dimensional gauge-theory computation to model.

\input{tex/summary.tex}

\input{tex/bloch.tex}
\input{tex/on_model.tex}
\input{tex/kt_flow.tex}
\input{tex/haldane.tex}
\input{tex/chair44.tex}

\section{Reproducing this document}
\label{sec:repro}

From the repository root:
\begin{verbatim}
make setup    # pip install -r tools/requirements.txt
make          # run all models, then build build/physics_models.pdf
make full     # longer Monte Carlo runs and one more ring size
\end{verbatim}
Each script writes figures to \texttt{build/figures/}, LaTeX fragments to \texttt{build/tex/} and raw numbers to \texttt{build/data/}.
The scripts check their own results (for example the AKLT ground-state energy and the Bloch correction coefficient) and stop the build if
a check fails. The Monte Carlo uses fixed seeds, so results are reproducible on the same software versions.

\section*{Limits of this document}
These computations are not part of the OpenAI release and are not reviewed by its authors. Statements attributed to the repository
come from its \texttt{CONTENTS.md} summaries and, for a few papers, their introductions; consult the papers themselves for exact hypotheses.
Numerical agreement with a formula does not establish a theorem, and where a comparison is only a shape comparison (for example the
second-moment mass against a transfer-matrix gap) the text says so.

\end{document}
